\scanpage{21}
peut néanmoins définir $\Phi_F$, ses sections sur l'ouvert $U\subset X$ sont les $A\subset U$ tels que $\forall x\in U$, $\exists$ voisinage (ouvert \unsure{si on veut}) $V_x$ \inserted{et $B_x\in F$} tels que $A\cap V_x=B_x\cap V_x$. Ici encore $\overline F=\Gamma(X,\Phi_F)$ \inserted{(plus généralement, $\Gamma(U,\Phi_F)=\overline F_U$ si $U$ ouvert de $X$)} satisfait a) b) c) d) (\unsure{peut-être|sans doute} pas e), ni $\overline F$ engendré par ses fermés\ldots), mais \unsure{$=$|même} si $X$ compact, on ne peut affirmer que \deleted{$\unsure{\Phi|F}$} $\overline F=F$. Par contre, si $F=\{\varnothing,X\}$, alors $\overline F$ est formé des parties à la fois ouvertes et fermées de $X$ (donc si $X$ est loc connexe, ce sont \unsure{les|des} réunions de comp. connexes).

\deleted{\unsure{plus général$^t$|plus généralement}}

Supposons p.ex. que $F=F_{\mathcal P}$, où $\mathcal P$ est décomposition cellulaire de $X$. Quand aura-t-on $\overline F=F$? Toute partie ouverte et fermée d'une cellule fermée est $\in\overline F$, il \uline{faut} donc que les cellules fermées soient connexes. Si on veut plus généralement que pour $\underline{\forall}$ ouvert $U\subset X$, $U\in F$, $\overline F_U=F_U$, on voit de \unsure{$=$|même} qu'il faut que $\underline{\forall}$ les cellules soient connexes (puisque elles sont fermées dans un ouvert $U\in F$). Cela est-il suffisant pour que $\overline F_U=F_U$ pour \unsure{tt} $U$ ouvert $\in F$?
\mspage{22}
